\begin{definition}[Strictly Increasing]
\label{def:strictly-increasing}
Let $A \subseteq \mathbb{R}$ and let $f : A \to \mathbb{R}$. We say that $f$ is \emph{strictly increasing} on $A$ if
\[
\forall x,y \in A\;
\bigl(
x < y \Rightarrow f(x) < f(y)
\bigr).
\]
\end{definition}

\begin{definition}[Strictly Decreasing]
\label{def:strictly-decreasing}
Let $A \subseteq \mathbb{R}$ and let $f : A \to \mathbb{R}$. We say that $f$ is \emph{strictly decreasing} on $A$ if
\[
\forall x,y \in A\;
\bigl(
x < y \Rightarrow f(x) > f(y)
\bigr).
\]
\end{definition}

\begin{definition}[Increasing at a Point]
\label{def:increasing-at-a-point}
Let $A \subseteq \mathbb{R}$, let $f : A \to \mathbb{R}$, and let $c \in A$. We say that $f$ is \emph{increasing at $c$} if there exists $\delta > 0$ such that for all
\[
x \in A \cap (c-\delta,c)
\quad \text{and} \quad
y \in A \cap (c,c+\delta),
\]
we have
\[
f(x) \leq f(c) \leq f(y).
\]
\end{definition}

\begin{definition}[Decreasing at a Point]
\label{def:decreasing-at-a-point}
Let $A \subseteq \mathbb{R}$, let $f : A \to \mathbb{R}$, and let $c \in A$. We say that $f$ is \emph{decreasing at $c$} if there exists $\delta > 0$ such that for all
\[
x \in A \cap (c-\delta,c)
\quad \text{and} \quad
y \in A \cap (c,c+\delta),
\]
we have
\[
f(x) \geq f(c) \geq f(y).
\]
\end{definition}












\begin{theorem}[Zero Derivative Implies Constant]
\label{thm:zero-derivative-implies-constant}
Let $I := [a,b] \subseteq \mathbb{R}$ be a closed interval, and let
$f : I \to \mathbb{R}$ be continuous on $I$ and differentiable on $(a,b)$.
If
\[
f'(x) = 0
\qquad \text{for all } x \in (a,b),
\]
then $f$ is constant on $I$.
\end{theorem}


\begin{corollary}[Equal Derivatives Imply Difference Is Constant]
\label{cor:equal-derivatives-constant-difference}
Let $I := [a,b] \subseteq \mathbb{R}$, and let $f,g : I \to \mathbb{R}$
be continuous on $I$ and differentiable on $(a,b)$. If
\[
f'(x) = g'(x)
\qquad \text{for all } x \in (a,b),
\]
then there exists $C \in \mathbb{R}$ such that
\[
f(x) = g(x) + C
\qquad \text{for all } x \in I.
\]
\end{corollary}



\begin{theorem}[Monotonicity via the Derivative]
\label{thm:monotonicity-derivative}
Let $I \subseteq \mathbb{R}$ be an interval, and let $f : I \to \mathbb{R}$
be differentiable on $I$. Then:
\begin{enumerate}[label=(\alph*)]
\item $f$ is increasing on $I$ if and only if
\[
f'(x) \geq 0
\qquad \text{for all } x \in I.
\]
\item $f$ is decreasing on $I$ if and only if
\[
f'(x) \leq 0
\qquad \text{for all } x \in I.
\]
\end{enumerate}
\end{theorem}



\begin{theorem}[First Derivative Test for Extrema]
\label{thm:first-derivative-test}
Let $I := [a,b] \subseteq \mathbb{R}$, let $f : I \to \mathbb{R}$ be continuous on $I$,
and let $c \in (a,b)$. Suppose $f$ is differentiable on $(a,c)$ and $(c,b)$.

\begin{enumerate}[label=(\alph*)]
\item If there exists $\delta > 0$ such that $(c-\delta,c+\delta) \subseteq I$ and
\[
f'(x) \geq 0 \quad \text{for } c-\delta < x < c,
\qquad
f'(x) \leq 0 \quad \text{for } c < x < c+\delta,
\]
then $f$ has a relative maximum at $c$.

\item If there exists $\delta > 0$ such that $(c-\delta,c+\delta) \subseteq I$ and
\[
f'(x) \leq 0 \quad \text{for } c-\delta < x < c,
\qquad
f'(x) \geq 0 \quad \text{for } c < x < c+\delta,
\]
then $f$ has a relative minimum at $c$.
\end{enumerate}
\end{theorem}


\begin{lemma}[Local Behavior from the Derivative]
\label{lem:local-derivative-behavior}
Let $I \subseteq \mathbb{R}$ be an interval, let $f : I \to \mathbb{R}$,
let $c \in I$, and suppose that $f$ is differentiable at $c$.

\begin{enumerate}[label=(\alph*)]
\item If $f'(c) > 0$, then there exists $\delta > 0$ such that
\[
f(x) > f(c)
\qquad \text{for all } x \in I \text{ with } c < x < c+\delta.
\]

\item If $f'(c) < 0$, then there exists $\delta > 0$ such that
\[
f(x) > f(c)
\qquad \text{for all } x \in I \text{ with } c-\delta < x < c.
\]
\end{enumerate}
\end{lemma}

\begin{theorem}[Darboux's Theorem]
\label{thm:darboux}
Let $I := [a,b] \subseteq \mathbb{R}$ and let $f : I \to \mathbb{R}$
be differentiable on $(a,b)$. If $k$ is a real number between $f'(a)$
and $f'(b)$, then there exists $c \in (a,b)$ such that
\[
f'(c) = k.
\]
\end{theorem}


















